\begin{slikaDesno}{fig/R_stick.pdf}
\PID 
На слици је приказан отпорни сензор у облику квадра, површине попречног пресека $S = 1\unit{mm^2}$, у укупне дужине 
$\ell = 20\unit{mm}$. 
Материјал од кога је отпорник израђен има Јунгов модул еластичности $E_{\rm Y} = 5\unit{MPa}$ и специфичну електричну
отпорност $\uprho = 50\unit{\Omega\,cm}$. Основе квадра су метализоване и ту су изведени прикључци на којима се мери његова отпорност $R$. 
Под дејством аксијалне силе $F$ отпорник се истеже, чиме се на механички начин мења 
његова отпорност. 
\end{slikaDesno}
\begin{enumerate}[label=(\alph*)]
    \item Сматрајући да се попречни пресек отпорника практично не мења, одредити зависност отпорности овог сензора 
    у зависности од примењене силе, $R = R(F)$. Израчунати силу пуне скале $F_{\rm FS}$, тако да се при тој сили 
    остварује максимално дозвољено истезање материјала $\Delta \ell = 5\% \, \ell$.

    \item Узимајући у обзир да се приликом истезања површина попречног пресека \textit{мења}, тако да запремина 
    отпорника остаје иста, релативну нелинеарност сензора ако је издужење пуне скале $\Delta \ell_{\rm FS} = 1\%\,\ell$.
\end{enumerate} 

\RESENJE

Јунгов модул еластичности даје везу између површинског напрезања (површинског напона) материјала, и његовог релативног истезања 
у задатом правцу\footnote{У општем случају, Јунгов модул је тензор, тако да омогућује да се материјал не истеже у по истом правцу 
у коме је напрезање. У овом задатку се то не разматра.}, односно 
\begin{equation}
    E_{\rm Y} = \dfrac{\upsigma}{\upepsilon} = \dfrac{F/S}{\Delta \ell/\ell} = \dfrac{F \ell}{S \Delta \ell},
    \label{eq:\ID:jung_uvod}
\end{equation}
Одавде се може добити резултат Хуковог закон, односно, коефицијент еластичности оваквог парчета материјала 
$F = \underbrace{\dfrac{S E_{\rm Y}}{\ell}}_k \Delta \ell$,
$k = 250\unit{mN/mm}$ . Отпорност приликом издужења 
од $\Delta \ell$ се онда може изразити у облику 
\begin{equation}
    R = \uprho \dfrac{\ell + \Delta \ell}{S}
      = \dfrac{\uprho \ell}{S} (1 + \upepsilon)
      = \underbrace{\uprho \dfrac{\ell}S}_{R_0} + \frac{\uprho \ell}{S^2 E_{\rm Y}} F 
      = 10\unit{k\Omega} + 2\unit{\dfrac{\Omega}{mN}} \, F
\end{equation}
Сила пуне скале се може израчунати помоћу израчунатог коефицијента еластичности, као 
$F_{\rm FS} = k {\Delta \ell_{\rm max}} = {5\%\,\ell}{k} = 500\unit{mN}$

(б) У случају да се деформација попречног пресека узима у обзир, то се може анализирати применом тоталног диференцијала. 
Запремина отпорника је $V = S\ell$. Налажењем тоталног диференцијала је 
$\de V = S \,\de \ell + \ell\, \de S$. Пошто се запремина не мења, то је $\de V = 0$, 
па се за промену дужине $\Delta \ell$ и промену површине $\Delta S$, може приближно 
писати $S\, \de \ell + \ell\, \de S = 0$, односно је 
$\Delta S = -\dfrac{\Delta \ell}{\ell} S = -\upepsilon S$. Заменом ових резултата у \ref{eq:\ID:jung_uvod} се онда добија 
\begin{equation}
 E_{\rm Y} = \dfrac{\dfrac{F}{S + \Delta S}}{\upepsilon} 
 = F\dfrac{\dfrac{1}{S - \upepsilon S}}{\upepsilon}
 = \dfrac{F}{S} \dfrac{1}{\upepsilon(1-\upepsilon)}
\end{equation}
На основу израза за отпорност је у овом случају 
\begin{equation}
    R = \uprho \dfrac{l + \Delta \ell}{S + \Delta S} = \dfrac{\uprho \ell}{S}\dfrac{1 + \upepsilon}{1 - \upepsilon}
     = \dfrac{\uprho \ell}{S} 
     \left(
        1 
        +
        \dfrac{2\upepsilon}{1 - \upepsilon}
     \right)
\end{equation}